% M20, 12.09.2026: formative/summative Evaluation zweckbezogen belegt;
% Forschungsablauf, Belegreichweiten und historische Standtrennung erhalten.
% M15, 11.09.2026: Absatzlogik und Lesbarkeit; methodische Einordnung,
% Literaturdarstellung, Belegarten und Standtrennung erhalten.
% ============================================================
% Kapitel 3 Auftakt + §3.1 Forschungsdesign — v03 (Claude)
%
% v03 (31.08., Anpassung an Blaupause-v02 = Zweiteilung 3.1+3.2):
%   · Auftakt kündigt nur noch 3.1 und 3.2 an; formativ/summativ-
%     Trennung ersetzt die 3.3-Ankündigung; ~80 Wörter (Budget 70–100).
%   · \ref{sec:durchfuehrung-untersuchung} → \ref{sec:iteratives-
%     vorgehen-evidenz}; \ref{sec:evaluationskonzept} ENTFERNT.
%   · Der reservierte Hevner-Guideline-3-Wortlaut (S. 83, "must be
%     rigorously demonstrated") wandert mit dem Evaluationsdesign
%     nach KAPITEL 7 (nicht mehr §3.3).
%   · TF-Tabelle bleibt 3-spaltig (Kapitel-Verweise in Spalte 3
%     enthalten) — funktional äquivalent zur 4-Spalten-Empfehlung
%     der v02-Blaupause §6.
%   · POSITIVE DREHUNG Absatz 3 (Autor-Auftrag): Einstieg jetzt
%     affirmativ ("Der tatsächliche Verlauf weist die DSRM-Aktivitäten
%     auf"), dann die zweifache Ordnungsrahmen-Rolle, dann die zwei
%     ehrlichen Sätze (kein Vorab-Protokoll; Mischlage beim Einstieg).
%     Inhaltlich identisch zur v02-Blaupause §2.2 — nur die Reihenfolge
%     führt mit der Stärke statt mit dem Vorbehalt.
%
% v02 (31.08., Review-Runde 1 — kritisch geprüft):
%   ÜBERNOMMEN:
%   ② Hevner-Präzisierung — PDF-Check: S. 82 trägt nur "must yield
%     utility … thorough evaluation"; das stärkere "must be rigorously
%     demonstrated" steht auf S. 83 (Guideline 3). Mein Satz hatte
%     S.-83-Inhalt mit S.-82-Zitat vermischt → S.-82-treue Fassung
%     übernommen; Guideline-3-Wortlaut (S. 83) RESERVIERT für §3.3.
%   ①-TEILWEISE: "ist eine …Studie" war stärker als das Kanon-Verb →
%     Absatz 1 jetzt "folgt einem …Vorgehen" (Kollegen-Fassung);
%     Auftakt neutraler ("Einordnung des Vorgehens in die DSR").
%   ZURÜCKGEWIESEN (Kanon-Konflikt): das Voll-Downgrade "orientiert
%     sich an der DSR" für Absatz 2 — die VERBINDLICHE Kapitel-3-
%     Blaupause (§1 Kernaussage + Sprachregel "wird als Design-Science-
%     Studie eingeordnet" = zulässig) verlangt die Klassifikation.
%     Absatz 2 behält "wird dem Design-Science-Paradigma zugeordnet".
%     Will der Autor das generelle Downgrade, ist ZUERST die Blaupause-
%     Kernaussage zu ändern.
%   AUFGELÖST (Runde 2, Autor+Kollege): Downgrade wird NICHT übernommen;
%     Endposition = "DSR als legitime Klassifikation des Vorgehens —
%     DSRM als nachträglich gewählter Ordnungsrahmen, nicht als
%     rückwirkend behauptetes Vorabprotokoll." Präzisierung übernommen:
%     "Das Vorgehen wird methodisch … zugeordnet" (Subjekt-Kontinuität
%     zu Absatz 1).
%   ③ (Chat-Kommentar, nicht Text): Klarstellung übernommen — Peffers
%     S. 56 belegt die Vereinbarkeit von Nichtlinearität/spätem
%     Einstieg mit der DSRM; die Quelle rechtfertigt NICHT die
%     Behauptung, DSRM sei vor ihrer bewussten Auswahl bereits
%     "angewandt" worden.
%
% Grundlage: Schreibblaupause-Kapitel-3 (chap3/, §4+§5) + die zwei
% lokalen Volltexte in chap3/3.1/. ALLE Fundstellen von Claude am PDF
% verifiziert (31.08.2026):
%   Hevner et al. (MISQ 28(1), 75–105):
%     S. 75  "seeks to extend the boundaries of human and organizational
%            capabilities by creating new and innovative artifacts"
%     S. 77  "creates and evaluates IT artifacts intended to solve
%            identified organizational problems"
%     S. 78  "build-and-evaluate loop is typically iterated a number of
%            times"; Ko-Evolution von Prozess und Artefakt
%     S. 79  Design als Suchprozess
%     S. 80  Framework Relevance/Rigor; Artefakttypen constructs/models/
%            methods/instantiations
%     S. 82  Guideline 1 (purposeful artifact) + "utility … thorough
%            evaluation" (Guideline 3)
%   Peffers et al. (JMIS 24(3), 45–77; DOI auf PDF-Titelseite):
%     S. 46  Aufzählung der sechs Aktivitäten
%     S. 52–56  Prozessbeschreibung im Detail
%     S. 56  "In reality, they may actually start at almost any step and
%            move outward"; "iterate back to activity 3"; nominelle
%            Struktur ≠ Prozesszwang
%
% Blaupausen-Konformität:
%   · Kapitelauftakt 1 Absatz (~100 Wörter), keine DSR-Definition,
%     keine Lücken-Wiederholung, kein KI-Absatz.
%   · 3.1 = 4 Absätze (Charakter → DSR/DSRM → ehrliche zeitliche
%     Einordnung → TF-Zuordnung) + ausgelagerte Tabelle.
%   · Einstiegspunkt: bewusst als Mischlage beschrieben (Blaupause §5
%     Absatz 3: nur benennen, wenn eindeutig — ist er nicht: praktisches
%     Zielbild + explorative Entwicklung).
%   · Sprachregeln eingehalten: "wird eingeordnet", "Ordnungsrahmen",
%     "rekonstruiert", KEIN "von Beginn an nach DSRM", KEIN
%     Wirksamkeits-Vorgriff.
%   · Systemneutralität: keine internen Namen (Core/Ledger/Steward);
%     MAF-Nennung zulässig (Kap. 2.4).
% Tabelle: chap3.1table.tex → tables/methodischer-zugang-teilfragen.tex.
% Bib: chap3.bib (hevner2004design, peffers2007design — verifiziert).
% MERKER: Verweis "Kapitel 1" für die Teilfragen ggf. durch \ref
%   ersetzen, sobald die Kapitel-1-Labels feststehen.
% MERKER: \ref{sec:iteratives-vorgehen-evidenz} = verbindliches Label
%   aus Blaupause-v02 — löst erst auf, wenn chap3.2 existiert (bis
%   dahin "??" im Kompilat, kein Fehler).
% ============================================================

\chapter{Forschungsmethodik}
\label{chap:forschungsmethodik}

Dieses Kapitel erläutert, wie die Arbeit aus der Entwicklung und
Untersuchung eines Agentensystems Erkenntnisse gewinnt.
Abschnitt~\ref{sec:forschungsdesign} begründet die methodische Einordnung;
Abschnitt~\ref{sec:iteratives-vorgehen-evidenz} beschreibt das iterative
Vorgehen und die verwendeten Belege. Unterschieden werden formative
Erkenntnisse, die während der Entwicklung den Entwurf beeinflussten,
und die abschließende Evaluation ausgewählter Systemeigenschaften in
Kapitel~\ref{chap:evaluation}.

\section{Forschungsdesign}
\label{sec:forschungsdesign}

Die Arbeit folgt einem explorativen, iterativen Entwicklungs- und
Evaluationsvorgehen. Untersucht wird, wie sich wiederkehrende
Projektkommunikation in nachvollziehbare Planungsartefakte überführen
lässt, die kontrolliert fortgeschrieben werden können. Im Mittelpunkt
steht das dafür entwickelte Artefakt: eine auf dem Microsoft Agent
Framework basierende Architektur einschließlich ihrer prototypischen
Realisierung. Die Untersuchung soll begründete Erkenntnisse darüber
liefern, wie agentische Verarbeitung mit festgelegten Prüf- und
Entscheidungsabläufen verbunden werden kann. Hinzu kommt eine empirische
Untersuchung ausgewählter Eigenschaften des Artefakts; eine allgemeine
Überlegenheit agentischer Systeme wird daraus nicht abgeleitet.

Das Vorgehen wird methodisch dem Design-Science-Paradigma zugeordnet.
Design-Science-Forschung gewinnt Erkenntnis, indem sie für ein identifiziertes
Problem neuartige Artefakte konstruiert und evaluiert
\cite[S.~75, 77]{hevner2004design}. Der Entwurf verläuft dabei als iterativer
Suchprozess, in dem sich Bauen und Bewerten in wiederholten Schleifen
abwechseln und sich Entwurfsprozess und Artefakt gemeinsam weiterentwickeln
\cite[S.~78--79]{hevner2004design}; da das Artefakt für das spezifizierte
Problem Nutzen erbringen soll, ist seine gründliche Evaluation zentral
\cite[S.~82]{hevner2004design}.

Zur Strukturierung des Forschungsprozesses
dient die Design Science Research Methodology (DSRM) nach Peffers et al. Sie
gliedert Design-Science-Vorhaben in sechs Aktivitäten -- Problemidentifikation
und Motivation, Definition der Lösungsziele, Design und Entwicklung,
Demonstration, Evaluation und Kommunikation
\cite[S.~46, 52--56]{peffers2007design} -- und ist ausdrücklich nicht als
zwingend lineare Abfolge angelegt: Forschungsvorhaben können an nahezu jedem
Punkt des Prozesses einsetzen, und von der Evaluation kann in Design und
Entwicklung zurückgekehrt werden \cite[S.~56]{peffers2007design}.

Der Verlauf dieser Arbeit umfasst Problemklärung, Formulierung von
Lösungszielen, iterative Gestaltung und formative Prüfung des Artefakts,
Demonstration sowie eine abschließende, begrenzte Evaluation. Die DSRM
dient als Ordnungsrahmen für die Darstellung dieser Aktivitäten.
Die Entwicklungszyklen werden anhand vorhandener Projekt- und
Laufzeitbelege rekonstruiert
(Abschnitt~\ref{sec:iteratives-vorgehen-evidenz} und
Kapitel~\ref{chap:explorative-problemklaerung}). Die technische Entwicklung
begann jedoch praxisorientiert, bevor diese methodische Einordnung
gewählt wurde. Es lag somit kein von Beginn an formalisiertes
DSRM-Protokoll vor. Der Projektbeginn verband ein praktisches Zielbild
mit explorativer Entwicklung und lässt sich keinem einzelnen der von
Peffers et al. beschriebenen Einstiegspunkte eindeutig zuordnen.

Formative Evaluation soll die Gestaltung eines Artefakts verbessern;
summative Evaluation soll seine Eignung für eine bestimmte Verwendung
beurteilen. Die Unterscheidung betrifft den Zweck der Bewertung und
nicht allein ihren Zeitpunkt \cite[S.~78--79]{venable2016feds}.

Die abschließende Evaluation in Kapitel~\ref{chap:evaluation}
untersucht ausgewählte Eigenschaften des Artefakts. Sie verbindet
referenzgestützte
Auswertungen dokumentierter Systemläufe, rückblickende Analysen
ausgewählter Entwicklungsfälle und ergänzende Kontrolltests, deren
Prüffälle vor der Ausführung spezifiziert wurden. Implementierungs-,
Eingabe- und Bewertungsstände sowie nachträgliche Revisionen werden
dort getrennt ausgewiesen. Die Nachweise stammen daher nicht sämtlich
aus einer gemeinsamen Ausführung an einem einheitlichen späteren
Endstand.

Die in Abschnitt~\ref{sec:einleitung-forschungsfragen} formulierten
Teilfragen werden mit
unterschiedlichen methodischen Zugängen bearbeitet.
Tabelle~\ref{t:methodischer-zugang-teilfragen} ordnet ihnen jeweils das
Erkenntnisziel, das verwendete Material und die Auswertungsform zu.
Die Art der Belege bestimmt dabei, welche Schlüsse aus den Ergebnissen
gezogen werden können.

% ============================================================
% Tabelle zu §3.1 — in Overleaf ablegen als:
%   tables/chap3/chap3.1table.tex
% Aufruf im Fließtext:
%   % ============================================================
% Tabelle zu §3.1 — in Overleaf ablegen als:
%   tables/chap3/chap3.1table.tex
% Aufruf im Fließtext:
%   % ============================================================
% Tabelle zu §3.1 — in Overleaf ablegen als:
%   tables/chap3/chap3.1table.tex
% Aufruf im Fließtext:
%   \input{tables/chap3/chap3.1table}
% Label: t:methodischer-zugang-teilfragen
% Inhalt = Blaupause-Kap-3 §5 Absatz 4 (TF-Zuordnung), sprachlich an
% die tatsächliche Evidenzlage angepasst (keine erfundenen Materialien).
% ============================================================

\begin{table}[tbp]
\centering

\caption[Methodische Zuordnung der Teilfragen]
{Zuordnung der Teilfragen zu Erkenntniszielen, Material und methodischem
Zugang. Eigene Darstellung.}
\label{t:methodischer-zugang-teilfragen}

\small
\renewcommand{\arraystretch}{1.15}

\begin{tabularx}{\textwidth}{
    >{\raggedright\arraybackslash}p{0.09\textwidth}
    >{\raggedright\arraybackslash}p{0.26\textwidth}
    >{\raggedright\arraybackslash}X
}
\toprule
\textbf{Teilfrage} &
\textbf{Erkenntnisziel} &
\textbf{Material und methodischer Zugang} \\
\midrule

TF1 &
Herausforderungen und Designanforderungen &
Nutzungsziel, ausgewählte frühe Entwicklungsstände, Läufe und
Entscheidungsnotizen;
formative Rekonstruktion und Design Rationale (Kapitel~\ref{chap:explorative-problemklaerung}) \\
\addlinespace[0.4em]

TF2 &
Gestaltung und Realisierung &
finaler Artefaktstand, Architektur und Frameworkabbildung;
Artefaktbeschreibung und technische Demonstration (Kapitel~\ref{chap:loesungsarchitektur} und~\ref{chap:maf-realisierung}) \\
\addlinespace[0.4em]

TF3 &
Eigenschaften und Trade-offs &
eingefrorene Evaluationsläufe, Referenzannotation und Challenge-Tests;
quantitative und qualitative Evaluation (Kapitel~\ref{chap:evaluation}) \\

\bottomrule
\end{tabularx}
\end{table}

% Label: t:methodischer-zugang-teilfragen
% Inhalt = Blaupause-Kap-3 §5 Absatz 4 (TF-Zuordnung), sprachlich an
% die tatsächliche Evidenzlage angepasst (keine erfundenen Materialien).
% ============================================================

\begin{table}[tbp]
\centering

\caption[Methodische Zuordnung der Teilfragen]
{Zuordnung der Teilfragen zu Erkenntniszielen, Material und methodischem
Zugang. Eigene Darstellung.}
\label{t:methodischer-zugang-teilfragen}

\small
\renewcommand{\arraystretch}{1.15}

\begin{tabularx}{\textwidth}{
    >{\raggedright\arraybackslash}p{0.09\textwidth}
    >{\raggedright\arraybackslash}p{0.26\textwidth}
    >{\raggedright\arraybackslash}X
}
\toprule
\textbf{Teilfrage} &
\textbf{Erkenntnisziel} &
\textbf{Material und methodischer Zugang} \\
\midrule

TF1 &
Herausforderungen und Designanforderungen &
Nutzungsziel, ausgewählte frühe Entwicklungsstände, Läufe und
Entscheidungsnotizen;
formative Rekonstruktion und Design Rationale (Kapitel~\ref{chap:explorative-problemklaerung}) \\
\addlinespace[0.4em]

TF2 &
Gestaltung und Realisierung &
finaler Artefaktstand, Architektur und Frameworkabbildung;
Artefaktbeschreibung und technische Demonstration (Kapitel~\ref{chap:loesungsarchitektur} und~\ref{chap:maf-realisierung}) \\
\addlinespace[0.4em]

TF3 &
Eigenschaften und Trade-offs &
eingefrorene Evaluationsläufe, Referenzannotation und Challenge-Tests;
quantitative und qualitative Evaluation (Kapitel~\ref{chap:evaluation}) \\

\bottomrule
\end{tabularx}
\end{table}

% Label: t:methodischer-zugang-teilfragen
% Inhalt = Blaupause-Kap-3 §5 Absatz 4 (TF-Zuordnung), sprachlich an
% die tatsächliche Evidenzlage angepasst (keine erfundenen Materialien).
% ============================================================

\begin{table}[tbp]
\centering

\caption[Methodische Zuordnung der Teilfragen]
{Zuordnung der Teilfragen zu Erkenntniszielen, Material und methodischem
Zugang. Eigene Darstellung.}
\label{t:methodischer-zugang-teilfragen}

\small
\renewcommand{\arraystretch}{1.15}

\begin{tabularx}{\textwidth}{
    >{\raggedright\arraybackslash}p{0.09\textwidth}
    >{\raggedright\arraybackslash}p{0.26\textwidth}
    >{\raggedright\arraybackslash}X
}
\toprule
\textbf{Teilfrage} &
\textbf{Erkenntnisziel} &
\textbf{Material und methodischer Zugang} \\
\midrule

TF1 &
Herausforderungen und Designanforderungen &
Nutzungsziel, ausgewählte frühe Entwicklungsstände, Läufe und
Entscheidungsnotizen;
formative Rekonstruktion und Design Rationale (Kapitel~\ref{chap:explorative-problemklaerung}) \\
\addlinespace[0.4em]

TF2 &
Gestaltung und Realisierung &
finaler Artefaktstand, Architektur und Frameworkabbildung;
Artefaktbeschreibung und technische Demonstration (Kapitel~\ref{chap:loesungsarchitektur} und~\ref{chap:maf-realisierung}) \\
\addlinespace[0.4em]

TF3 &
Eigenschaften und Trade-offs &
eingefrorene Evaluationsläufe, Referenzannotation und Challenge-Tests;
quantitative und qualitative Evaluation (Kapitel~\ref{chap:evaluation}) \\

\bottomrule
\end{tabularx}
\end{table}
